\documentclass[UTF8,12pt,a4paper]{ctexart}

\usepackage{amsmath,amssymb,bm}
\usepackage{booktabs}
\usepackage{geometry}
\usepackage{longtable}
\usepackage{array}
\usepackage{hyperref}
\usepackage{graphicx}
\usepackage{tikz}
\usetikzlibrary{arrows.meta,positioning}

\geometry{left=2.6cm,right=2.6cm,top=2.5cm,bottom=2.5cm}
\hypersetup{colorlinks=true,linkcolor=blue,citecolor=blue,urlcolor=blue}
\newcommand{\Rcal}{\mathcal{R}}
\newcommand{\Fsen}{\mathcal{F}}
\title{基于支撑域约束低速模板与结构化频率模型的叶尖间隙--振动联合辨识理论}
\author{}
\date{}

\begin{document}

\maketitle

\begin{abstract}
径向间隙变化、低速模板离散化误差和有限传感器布局均会改变高速完整波形中可用于区分振动阶次的信息。本文建立一套与冻结仿真程序一致的联合辨识理论：噪声先加到完整数据采集（data acquisition，DAQ）时间记录，再经一次转速参考（once-per-revolution，OPR）映射形成空间波形；低速侧在实测支撑域内采用按完整转次分组交叉验证的自适应 Savitzky--Golay（SG）平滑，并以分段三次 Hermite 插值（PCHIP）形成连续模板；高速侧采用相对于低速状态的间隙增量前向模型，并分别处理单同步、单异步、双同步和同步--异步四类已知频率结构。对于双同步模型，Funnel V2 使全部离散 EO 对先经历一次间隙、空间偏移、幅值和相位的联合 Gauss--Newton 更新，再对 Top--24 二次更新，并仅对 Top--3 执行完整非线性精修。该流程的职责是提高真 EO 候选召回率，而不是保证近退化 EO 模型在最终残差竞争中必胜。本文进一步将结构可辨识性、给定离散结构后的连续参数条件可恢复性、有限噪声下的实用可辨识性和离散 EO 模型选择分层讨论，并给出支撑域覆盖、模板导数质量、有效幅噪比、传感器独立相位信息、候选子空间分离度、优化收敛和最终残差裕量等可检验条件。仿真表明，在限定的三传感器、8 转、10--20 dB DAQ 信噪比和分量幅值不低于 $0.1\,\mathrm{mm}$ 的代表性范围内，四类结构化模型共 42 组测试均满足预定参数误差判据；同时，近共线 EO 子空间与确定性模板偏差可导致离散模型选择翻转，因此该结论不外推为任意 EO 组合下的唯一可辨识保证。
\end{abstract}

\noindent\textbf{关键词：}叶尖间隙；叶片振动；完整波形；低速模板；结构化频率；可辨识性；变量投影

\section{引言}

叶尖间隙是反映旋转机械气动性能和运行安全裕度的重要参数。非接触式叶尖传感器能够记录叶尖经过探头敏感区域时的高速电压波形\cite{Yu2020,Muller1997}。与主要利用叶尖到达时刻或其相位信息的方法\cite{Diamond2021}相比，完整波形还包含幅值和局部形状信息，因此可为叶尖间隙和叶片振动的联合估计提供观测基础\cite{Heller2020}。

传统叶尖定时通常利用阈值交点、峰值位置或其他脉冲特征确定叶片到达时刻，并将其相对于无振动参考时刻的偏差换算为叶尖周向振动位移。径向间隙变化会改变传感器脉冲的幅值、宽度和局部形状，即使叶片真实周向位置不变，也可能使提取的到达时刻发生偏移，并将该偏差传递到后续振动参数识别中。

完整波形方法通过低速基准脉冲与高速实测脉冲之间的位置差异识别振动。高速径向间隙发生变化时，高速脉冲在发生周向位置变化的同时还会产生幅值和形状变化，因而不再是低速模板的单纯平移。若模型只允许低速模板发生位置移动，间隙变化产生的部分波形差异就可能被解释为周向振动位移，从而影响振动幅值和相位的识别。

为此，本文利用静态多间隙数据描述径向间隙变化引起的脉冲幅值和形状变化，利用低速旋转数据保存当前叶片与传感器组合的实测基准波形。低速模板与静态响应面共同确定局部间隙路径；高速阶段保持该局部关系不变，只估计相对于低速状态的径向间隙增量，并与周向振动参数联合求解。由此避免仅通过移动低速模板解释高速间隙变化产生的波形差异。

本文首先从叶尖定时和完整波形匹配两方面说明径向间隙变化对振动识别的影响，并给出基本建模假设。随后说明低速旋转数据和静态多间隙数据的功能分工，建立径向间隙增量与周向振动位移共同作用的完整波形模型，最后给出候选生成、投影局部搜索和完整非线性优化流程。

\section{研究问题与建模假设}

\subsection{叶尖间隙--振动联合识别问题}

本文研究的问题是：当高速运行状态的径向间隙不同于低速基准状态时，如何利用同一组叶尖传感器完整波形，同时识别相对径向间隙变化和叶尖周向振动，并避免将间隙变化引起的波形差异误认为振动位移。为明确该问题，设 $g$ 为叶尖与传感器之间沿敏感方向的等效间隙，在本文试验布置下近似对应径向间隙；设 $x$ 为叶尖经过传感器敏感区域时的局部周向坐标，$u(t)$ 为叶尖振动在周向上的等效位移。

传统叶尖定时根据阈值交点、峰值或相关定位等脉冲特征确定叶片到达时刻。设叶片周向振动对应的真实到达时刻偏差为 $\Delta t_u$，径向间隙变化引起的特征定位偏差为 $\delta t_g$，其他测量误差为 $\varepsilon_t$，则提取的到达时刻偏差可表示为
\begin{equation}
\widehat{\Delta t}=\Delta t_u+\delta t_g+\varepsilon_t.
\label{eq:btt_time_error}
\end{equation}
在局部周向速度 $V_{\mathrm{tip}}$ 已知的条件下，传统叶尖定时给出的周向位移为
\begin{equation}
\widehat u_{\mathrm{BTT}}
=V_{\mathrm{tip}}\widehat{\Delta t}
\approx u+V_{\mathrm{tip}}\delta t_g+V_{\mathrm{tip}}\varepsilon_t,
\label{eq:btt_displacement}
\end{equation}
其中，$u\approx V_{\mathrm{tip}}\Delta t_u$ 为真实周向振动位移。式\eqref{eq:btt_displacement}表明，径向间隙变化一旦使脉冲幅值、宽度或局部轮廓发生变化，并进一步改变特征定位结果，其影响就会以附加周向位移的形式进入振动识别。

完整波形方法不再依赖单一特征点，而是通过低速基准模板与高速实测脉冲的整体匹配估计周向位移，但间隙与振动的混淆并不会因此自动消失。为说明这一点，在单传感器、单脉冲和小扰动条件下，记低速基准模板为 $T(x)$，模板的空间导数为 $T_x(x)$。若无振动波形关于间隙和周向位置的局部响应记为 $v(g,x)$，则低速基准间隙路径 $g^{\mathrm{low}}(x)$ 附近的间隙灵敏度定义为
\begin{equation}
S_g(x)=
\left.
\frac{\partial v(g,x)}{\partial g}
\right|_{g=g^{\mathrm{low}}(x)}.
\label{eq:problem_gap_sensitivity}
\end{equation}
在当前脉冲窗口内，将周向位移 $u$ 和径向间隙增量 $\Delta g$ 视为给定的小量，并忽略二者的乘积及其他二阶项，则高速波形相对于低速模板的一阶变化可写为
\begin{equation}
v^{\mathrm{high}}(x)-T(x)
\approx -T_x(x)u+S_g(x)\Delta g.
\label{eq:problem_linearization}
\end{equation}
式\eqref{eq:problem_linearization}中，第一项表示周向位移引起的波形平移，第二项表示径向间隙变化引起的幅值和形状变化。若识别模型只允许固定模板发生周向平移，即用 $-T_x(x)\widetilde u$ 解释全部波形变化，则在当前窗口内有近似关系
\begin{equation}
\widetilde u-u
\approx
-\Delta g\,
\frac{\langle T_x,S_g\rangle}
{\langle T_x,T_x\rangle},
\label{eq:gap_vibration_bias}
\end{equation}
其中，$\langle\cdot,\cdot\rangle$ 表示在有效脉冲窗口内的离散内积，$\widetilde u$ 为固定模板模型估计的周向位移。只要间隙灵敏度 $S_g$ 在模板平移方向 $T_x$ 上具有非零投影，固定模板匹配就会把部分间隙响应分配给周向位移。本文需要解决的正是这种由波形灵敏度重叠造成的间隙--振动混淆，而不是单纯追求更小的电压拟合残差。式\eqref{eq:problem_linearization}和式\eqref{eq:gap_vibration_bias}仅用于说明研究问题，后续识别采用完整非线性波形模型。

因此，在静态多间隙标定数据、低速基准波形和高速完整波形均可获得的条件下，本文估计各有效传感器相对于低速状态的等效径向间隙增量，以及由多传感器共同约束的周向振动频率、幅值和参考相位。该问题属于低速标定约束下的局部联合识别，所得间隙量不解释为缺少独立间隙基准条件下的绝对间隙测量值。

\subsection{基本建模假设}

为使低速模板、静态间隙响应面和高速观测数据能够在同一前向模型中建立联系，本文采用以下基本假设：
\begin{enumerate}
  \item 低速工况下，目标频带内的叶片振动相对于目标识别精度可忽略，或其在多转聚合后形成的系统性分量可忽略；同一目标叶片的重复脉冲经统一坐标转换后具有良好重复性。因此，多转聚合波形可作为当前叶片--传感器组合的低速基准模板。
  \item 在所选分析窗口内，各传感器对应测点处的平均径向间隙变化时间尺度远大于目标振动周期。因此，每个测点相对于低速状态的径向变化可用窗口内恒定的间隙增量 $\Delta g_s$ 表示。
  \item 在本文研究工况和目标频带内，叶片振动对测量波形的主要动态作用表示为叶尖周向等效位移，高速与低速之间沿径向的主要差异表示为准静态间隙增量。由此，动态位移采用周向分量描述，径向方向用于描述工况之间的准静态间隙差异。
  \item 在单个传感器的有效敏感窗口内，叶尖局部几何与传感器安装关系从低速到高速保持足够稳定。因此，由低速模板与静态响应面联合标定得到的等效局部间隙路径可用于描述高速状态下周向位置变化所对应的局部等效间隙变化。
  \item 经信号极性、基线、电压尺度和空间坐标校正后，静态标定、低速旋转和高速旋转条件下的局部间隙--电压响应具有可比性。因此，静态响应面可用于计算高速状态相对于低速状态的间隙增量响应。
  \item 低速和高速噪声分别在各自完整 DAQ 时间记录上生成，噪声加入后才进行 OPR 映射和空间截取；因此，空间辨识窗口中的实际信噪比由映射、窗口及波形局部幅值共同决定，不与名义 DAQ 信噪比混为一谈。
  \item 频率结构类别在一次求解中预先给定为单同步、单异步、双同步或同步--异步之一，但结构内的 EO 或连续频率未知。结构类别已知不等于具体频率已知，也不构成真实 EO 先验。
  \item 所有参与拟合的高速样本均满足支撑域契约：在允许的公共偏移、振动幅值和间隙范围内，其全部可能查询位置仍落在低速模板与静态响应面的共同安全支撑域内。算法不以插值外推补偿缺失物理信息。
\end{enumerate}

上述假设只定义前向模型的适用域，并不自动给出唯一可辨识性。尤其是，即使前向模型在真参数处精确，有限传感器布局下不同 EO 模型的局部灵敏度子空间仍可能近共线；反之，某一噪声实现下选择错误 EO 也不能直接推出结构不可辨识。第\ref{sec:identifiability}节将这两类问题分层说明。

\section{联合识别方法}

\subsection{方法总体框架}

本文方法按照信息进入模型的顺序分为五个阶段，如图\ref{fig:method_framework}所示。第一阶段在完整低速和高速 DAQ 时间记录上分别加入规定噪声，并记录名义与实际 DAQ 信噪比。第二阶段利用 OPR 将时间采样映射到局部周向坐标，同时计算低速观测支撑、高速观测支撑和静态响应面支撑。第三阶段在低速公共支撑域内进行空间分箱，并按完整转次分组交叉验证选择 SG 窗宽，再通过 PCHIP 得到连续实测模板及稳定导数。第四阶段将低速模板与静态响应面进行有界匹配，确定低速等效基准间隙和配准量，并建立只描述相对间隙变化的增量前向模型。第五阶段在预先给定的四类频率结构之一内生成候选，并在完整非线性电压模型中联合估计高速间隙增量、公共空间偏移和振动参数。

静态间隙库描述径向间隙变化引起的波形变化，低速模板保存当前叶片与传感器组合的实际基准波形。二者通过低速等效标定建立联系：响应面不替换实测模板，在高速阶段仅用于计算相对于低速状态的间隙增量响应。候选生成、候选召回和最终模型选择是三个不同环节；任何一级近似得分都不作为最终 EO 正确性的理论保证。

\begin{figure}[htbp]
\centering
\resizebox{\textwidth}{!}{%
\begin{tikzpicture}[
  >=Latex,
  data/.style={draw,rounded corners,fill=gray!8,text width=33mm,minimum height=10mm,align=center},
  proc/.style={draw,rounded corners,text width=37mm,minimum height=11mm,align=center},
  resultbox/.style={draw,rounded corners,fill=gray!8,text width=39mm,minimum height=11mm,align=center},
  flow/.style={->,thick}
]
\node[data] (static) at (0,0) {静态多间隙数据\\$\{g_i,v_i^{\mathrm{stat}}\}$};
\node[data] (low) at (5.1,0) {低速旋转数据\\$v_{s\ell}^{\mathrm{low}}(t)$};
\node[data] (high) at (10.2,0) {高速旋转数据\\$v_s^{\mathrm{high}}(t)$};

\node[proc] (surface) at (0,-2.3) {空间坐标统一\\构建 $\Rcal(g,x)$};
\node[proc] (template) at (5.1,-2.3) {时间--空间转换与公共网格重采样\\构建 $T_s^{\mathrm{low}}(x)$};
\node[proc] (hsample) at (10.2,-2.3) {时间--空间转换\\形成 $\{t_{sn},x_{sn},v_{sn}\}$};

\node[proc] (cal) at (3.4,-4.7) {低速等效标定\\确定局部间隙路径与配准参数};
\node[proc] (joint) at (9.0,-4.7) {固定低速标定关系\\完整波形联合识别};
\node[resultbox] (result) at (9.0,-6.7) {相对径向间隙变化\\振动频率、幅值和相位};

\draw[flow] (static) -- (surface);
\draw[flow] (low) -- (template);
\draw[flow] (high) -- (hsample);
\draw[flow] (surface) |- (cal.west);
\draw[flow] (template) -- (cal.north);
\draw[flow] (cal) -- (joint);
\draw[flow] (hsample) -- (joint.north);
\draw[flow] (joint) -- (result);
\end{tikzpicture}%
}
\caption{叶尖间隙与振动联合识别方法总体流程}
\label{fig:method_framework}
\end{figure}

\subsection{时间--空间坐标转换与低速模板构建}

叶尖传感器首先在时间域记录电压序列，而静态响应面和完整波形模型均以叶尖相对于传感器的局部周向位置为自变量。因此，在构建低速模板和高速观测模型之前，需要先规定采集层噪声协议，再利用转角参考将各转次的采样时刻转换为统一空间坐标。
仿真中的噪声协议在坐标变换之前定义。对低速或高速的完整无噪声 DAQ 电压记录 $\bm v_{\mathrm{DAQ}}^0$，定义交流有效值、噪声标准差和加噪记录为
\begin{equation}
s_{\mathrm{DAQ}}
=\operatorname{RMS}\!\left(\bm v_{\mathrm{DAQ}}^0-\overline v_{\mathrm{DAQ}}^0\bm 1\right),
\quad
\sigma_{\mathrm{DAQ}}=s_{\mathrm{DAQ}}10^{-\mathrm{SNR}_{\mathrm{DAQ}}/20},
\quad
\bm v_{\mathrm{DAQ}}=\bm v_{\mathrm{DAQ}}^0+\sigma_{\mathrm{DAQ}}\bm\varepsilon,
\label{eq:daq_noise}
\end{equation}
其中 $\bm\varepsilon\sim\mathcal N(\bm0,\bm I)$。低速和高速记录分别用各自完整无噪声记录的 $s_{\mathrm{DAQ}}$ 设置噪声；不对每个随机种子的噪声实现二次缩放到精确名义信噪比。完成式\eqref{eq:daq_noise}后，才对加噪记录执行 OPR 映射、可信空间区间截取和辨识。因此，名义 $\mathrm{SNR}_{\mathrm{DAQ}}$ 是采集层协议，可信空间窗口内的局部电压信噪比或振动增量信噪比是结果量，二者不能互换。

设相邻两个 OPR 时刻为 $t_m^{\mathrm{OPR}}$ 和 $t_{m+1}^{\mathrm{OPR}}$，两参考位置之间的已知转角增量为 $\Delta\theta_m^{\mathrm{OPR}}$。在该时间区间内，累计转角定义为
\begin{equation}
\theta(t)=\theta_m^{\mathrm{OPR}}
+\Delta\theta_m^{\mathrm{OPR}}
\frac{t-t_m^{\mathrm{OPR}}}
{t_{m+1}^{\mathrm{OPR}}-t_m^{\mathrm{OPR}}},
\qquad
t\in[t_m^{\mathrm{OPR}},t_{m+1}^{\mathrm{OPR}}],
\label{eq:time_to_angle}
\end{equation}
其中，$\theta_m^{\mathrm{OPR}}$ 为第 $m$ 个参考位置对应的累计转角；单槽 OPR 每转提供一次参考时，$\Delta\theta_m^{\mathrm{OPR}}=2\pi$，多槽或不等槽 OPR 则采用相邻参考位置之间的已知槽角。该映射利用实测 OPR 时刻反映不同转次或相邻槽区间的平均转速变化，并在相邻参考时刻之间采用角速度近似恒定的分段线性表示。

对第 $s$ 个传感器的第 $\ell$ 次目标叶片经过事件，记其名义参考转角为 $\theta_{s\ell}^{\mathrm{ref}}$，叶尖参考半径为 $R_{\mathrm{tip}}$。采样时刻 $t_{s\ell n}$ 对应的局部周向空间坐标定义为
\begin{equation}
x_{s\ell n}
=R_{\mathrm{tip}}
\left[\theta(t_{s\ell n})-\theta_{s\ell}^{\mathrm{ref}}\right].
\label{eq:local_space_coordinate}
\end{equation}
其中，$n$ 为单次叶片经过窗口内的采样点编号，$x=0$ 对应名义参考位置，正方向由转子旋转方向统一规定。式\eqref{eq:local_space_coordinate}把不同转速、不同转次下的时间采样点转换为具有相同物理含义的局部周向位置；若采用局部角坐标而非弧长坐标，则省略比例因子 $R_{\mathrm{tip}}$，但全文必须保持同一坐标单位。各传感器的局部周向坐标采用统一的物理单位、正方向和尺度，使后续共享的周向振动位移 $u(t)$ 在不同传感器模型中具有一致的物理含义。

设第 $s$ 个传感器在第 $\ell$ 个低速转次记录的原始电压为 $v_{s\ell}^{\mathrm{low}}(t_{s\ell n})$。将离散点按间距 $\Delta x_b$ 投入公共空间网格 $\{x_p\}_{p=1}^{N_x}$，每个分箱仅使用落入该箱且满足最小样本数要求的点，得到分箱均值 $\bar v_{s}^{\mathrm{bin}}(x_p)$。为避免同一转次同时进入训练集和验证集而低估泛化误差，以完整转次为分组单位进行 $K$ 折交叉验证。对候选物理窗宽 $w\in\mathcal W$，记相应奇数 SG 跨度为 $m(w)$，则选择
\begin{equation}
\widehat w_s
=\arg\min_{w\in\mathcal W}
\frac{1}{K}\sum_{k=1}^{K}
\operatorname{RMSE}\!\left(
\mathcal P_{w}^{(-k)},\mathcal D_{s}^{(k)}
\right),
\label{eq:spatial_resampling}
\end{equation}
其中 $\mathcal D_s^{(k)}$ 为第 $k$ 个完整转次组，$\mathcal P_w^{(-k)}$ 为其余转次经分箱和窗宽 $w$ 的 SG 平滑后得到的预测模板。冻结程序采用最小分组预测误差规则选择 $\widehat w_s$，而不是利用高速辨识结果反向选择窗宽。对全部低速转次重新分箱并以 $\widehat w_s$ 平滑后，记离散结果为 $\widetilde T_s(x_p)$；连续低速模板定义为
\begin{equation}
T_s^{\mathrm{low}}(x)
:=\operatorname{PCHIP}\!\left(\{x_p,\widetilde T_s(x_p)\}_{p=1}^{N_x};x\right).
\label{eq:low_template}
\end{equation}
PCHIP 只把离散平滑模板变为形状保持的连续查询函数，不恢复分箱已丢失的信息。候选生成所需的空间导数由 SG 平滑后的离散模板求导并再插值，避免直接对含噪硬分箱序列差分。$T_s^{\mathrm{low}}(x)$ 保存目标叶片、传感器安装、局部几何和测量链路共同形成的实际低速基准波形；其函数误差和导数误差必须分别评价。

高速数据采用相同的式\eqref{eq:time_to_angle}和式\eqref{eq:local_space_coordinate}转换。后文分别以 $t_{sn}$、$x_{sn}$ 和 $v_{sn}$ 表示第 $s$ 个传感器第 $n$ 个高速观测点的采样时刻、局部周向坐标和实测电压。其中，$x_{sn}$ 用于查询空间波形，$t_{sn}$ 保留用于计算振动相位。由此，波形形状在空间域中比较，振动随时间的演化仍由原始采样时刻描述。

可信辨识窗口不是单独经验截取，而由支撑域契约确定。记低速模板观测支撑、静态响应面空间支撑和数值算子边界裕量分别为 $\mathcal X_{L,s}$、$\mathcal X_R$ 和 $m_{\mathrm{num}}$，则共同安全支撑为
\begin{equation}
\mathcal X_{\mathrm{safe},s}
=\left(\mathcal X_{L,s}\cap\mathcal X_R\right)\ominus m_{\mathrm{num}}.
\label{eq:safe_support}
\end{equation}
若公共偏移满足 $d\in[d_{\min},d_{\max}]$，结构内各振动分量幅值上界为 $A_q^{\max}$，则 $u_{\max}=\sum_q A_q^{\max}$，可进入辨识的高速名义坐标必须满足
\begin{equation}
\mathcal X_{H,s}^{\mathrm{cert}}
=\left[
x_{mathrm{safe},s}^{-}+d_{\max}+u_{\max},
x_{mathrm{safe},s}^{+}+d_{\min}-u_{\max}
\right].
\label{eq:certified_high_support}
\end{equation}
于是对任意允许参数都有 $x_{sn}-d-u(t_{sn})\in\mathcal X_{\mathrm{safe},s}$。只有 $x_{sn}\in\mathcal X_{H,s}^{\mathrm{cert}}$ 的样本进入辨识；若该区间为空或规定的间隙域超出静态间隙库支撑，则停止求解，而不是外推。OPR 映射决定观测点落在何处，低速/静态支撑决定模型能够在哪些位置被查询，振幅和偏移上界进一步把该域收缩为真正的辨识窗口。

\subsection{静态间隙响应面与低速等效标定}

\subsubsection{间隙方向基函数与静态响应面构造}

静态标定用于建立无振动条件下电压关于径向间隙 $g$ 和局部周向坐标 $x$ 的响应关系。本文采用一个共享响应面作为统一间隙库，并假定不同传感器相对于该间隙库的主要差异可由后续的电压增益、偏置及空间配准参数表示；该假定应通过各传感器低速模板的重构误差检验。对固定的周向位置，本文采用常数项、反比项和对数项组成的低阶间隙基函数
\begin{equation}
\phi_0(g)=1,
\qquad
\phi_1(g)=\frac{1}{g},
\qquad
\phi_2(g)=\ln\left(\frac{g}{g_{\mathrm{ref}}}\right).
\label{eq:gap_basis}
\end{equation}
式中，$g_{\mathrm{ref}}>0$ 为参考间隙，用于使对数函数的自变量无量纲；$\phi_0$、$\phi_1$ 和 $\phi_2$ 分别为常数、反比和对数基函数。常数项给出间隙方向拟合的截距，反比项描述间隙减小时传感器响应增强的主要趋势，对数项修正有限标定范围内单一反比关系不能覆盖的缓变非线性。上述基函数是面向标定区间的经验表示，不作为严格的电场解析解。基于式\eqref{eq:gap_basis}，径向间隙--周向位置响应面定义为
\begin{equation}
\Rcal(g,x)
=\beta_0(x)+\frac{\beta_1(x)}{g}
+\beta_2(x)\ln\left(\frac{g}{g_{\mathrm{ref}}}\right),
\label{eq:response_surface}
\end{equation}
式中，$\Rcal(g,x)$ 为静态标定响应面；$\beta_0(x)$、$\beta_1(x)$ 和 $\beta_2(x)$ 为三个间隙基函数在周向位置 $x$ 处的系数，而不是预先规定形状的空间函数。允许这些系数随 $x$ 变化，是为了描述叶尖进入、正对和离开传感器敏感区域时，拟合截距、间隙灵敏度和非线性程度随周向位置发生的变化。

设静态标定包含 $N_g$ 个已知正间隙 $g_i$ 和 $N_x$ 个空间采样点 $x_p$。在固定空间位置 $x_p$ 处，将不同间隙下的实测电压和相应基函数值分别组成观测向量与标定矩阵：
\begin{equation}
\bm y_p=
\begin{bmatrix}
v_{1p}^{\mathrm{stat}} & \cdots & v_{N_gp}^{\mathrm{stat}}
\end{bmatrix}^{\mathrm T},
\qquad
\bm B_g=
\begin{bmatrix}
1 & 1/g_1 & \ln(g_1/g_{\mathrm{ref}})\\
\vdots & \vdots & \vdots\\
1 & 1/g_{N_g} & \ln(g_{N_g}/g_{\mathrm{ref}})
\end{bmatrix}.
\label{eq:cal_matrix}
\end{equation}
式中，$v_{ip}^{\mathrm{stat}}$ 为已知间隙 $g_i$、空间位置 $x_p$ 处的静态实测电压；$\bm y_p\in\mathbb R^{N_g}$ 为该位置的观测向量；$\bm B_g\in\mathbb R^{N_g\times3}$ 为间隙方向标定矩阵。三个系数可独立估计要求 $N_g\geq3$ 且 $\bm B_g$ 具有数值满列秩；若采用留一间隙重构检验，则应有 $N_g\geq4$。定义位置 $x_p$ 处的响应面系数向量为 $\bm\beta_p=[\beta_0(x_p),\beta_1(x_p),\beta_2(x_p)]^{\mathrm T}$，其最小二乘估计为
\begin{equation}
\bm\beta_p^*
=\arg\min_{\bm\beta_p}
\left\lVert \bm B_g\bm\beta_p-\bm y_p\right\rVert_2^2
\label{eq:cal_ls}
\end{equation}
式中，$\bm\beta_p^*$ 为 $\bm\beta_p$ 的估计值。具体而言，先在每个离散周向位置 $x_p$ 利用多个已知间隙估计三项系数，再沿周向位置对系数序列进行平滑插值，从而得到连续的 $\beta_j(x)$、响应面 $\Rcal(g,x)$ 及其间隙和空间导数。基函数形式及平滑结果通过留一间隙重构误差、$\partial\Rcal/\partial g$ 的符号与连续性以及标定矩阵的数值条件进行检验；后续识别仅在静态间隙库具有稳定覆盖的范围内进行。

\subsubsection{低速模板在静态响应面中的有界等效标定}

静态响应面提供统一的径向间隙--周向位置响应关系，不同传感器的安装零位、电压尺度和空间尺度通过传感器相关参数进行校正。为确定当前传感器在响应面中的低速查询路径，同时保留其实际基准波形，需要将实测模板 $T_s^{\mathrm{low}}(x)$ 与响应面进行有界匹配。

低速等效间隙路径和局部坐标到静态间隙库坐标的映射分别写为
\begin{equation}
g_s^{\mathrm{low}}(x)
=g_s^{(0)}+\mu_s^{\mathrm{cal}}(x-\tau_s),
\qquad
x_s^{\mathrm{lib}}(x)=\zeta_s(x-\tau_s).
\label{eq:low_cal_path}
\end{equation}
其中，$g_s^{(0)}$ 为低速模型等效基准径向间隙，$\mu_s^{\mathrm{cal}}$ 为径向间隙随局部周向位置变化的等效梯度，$\tau_s$ 为模板与间隙库的周向中心配准参数，$\zeta_s$ 为周向尺度修正系数。由此将低速查询路径明确写为静态响应面中的参数曲线
\begin{equation}
\Gamma_s(x)=
\left(g_s^{\mathrm{low}}(x),x_s^{\mathrm{lib}}(x)\right).
\label{eq:query_path}
\end{equation}
$\Gamma_s(x)$ 描述单个传感器有效窗口内，等效径向间隙随局部周向位置变化的标定路径。

为使响应面路径在有效标定域内逼近低速实测模板，定义路径参数向量 $\bm\vartheta_s=[g_s^{(0)},\mu_s^{\mathrm{cal}},\tau_s,\zeta_s]^{\mathrm T}$，并求解以下有界标定问题：
\begin{equation}
(\widehat{\bm\vartheta}_s,\widehat\kappa_s,\widehat b_s)
=\arg\min_{\substack{\bm\vartheta_s\in\Omega_s,\,\kappa_s>0\\b_s\in\mathbb R}}
\left\{
\sum_p
\left[
T_s^{\mathrm{low}}(x_p)-b_s
-\kappa_s\Rcal\!\left(
g_s^{\mathrm{low}}(x_p),x_s^{\mathrm{lib}}(x_p)
\right)
\right]^2
+P_{\mathrm{cal,out}}
\right\}
\label{eq:low_calibration}
\end{equation}
式中，$\Omega_s$ 为由静态间隙范围、有效空间窗口和物理可行范围共同确定的有界参数域，并在统一周向坐标正方向后约束 $\zeta_s>0$；$\kappa_s$ 和 $b_s$ 分别为响应面到第 $s$ 个传感器低速模板的电压增益和偏置；$P_{\mathrm{cal,out}}$ 为查询路径越出响应面支持范围时的惩罚项。统一电压极性后取 $\kappa_s>0$，以避免增益符号与响应方向相互补偿。对给定的 $\bm\vartheta_s$，$\kappa_s$ 和 $b_s$ 可由带上述符号约束的线性最小二乘确定，外层只需搜索查询路径参数。$b_s$ 仅为低速标定辅助参数，不进入高速主识别；$\kappa_s$ 用于将统一间隙库的电压变化尺度映射到当前传感器。

式\eqref{eq:low_calibration}用于确定低速模板在静态响应面中的等效查询路径，而不是用响应面拟合值替换实测模板。所得参数表示当前叶片与传感器组合在有效窗口内的等效标定关系。

低速标定结果通过增益方向、响应面覆盖范围以及重复数据和不同初值下查询路径的一致性进行检查。模板拟合残差与标定路径稳定性共同作为标定结果进入高速识别的依据。

\subsubsection{基于实测模板的间隙增量模型}

完成低速等效标定后，低速基准间隙和高速查询间隙分别定义为
\begin{align}
g_s^{\mathrm{base}}(x)
&=g_s^{(0)}+\mu_s^{\mathrm{cal}}(x-\tau_s),
\label{eq:g_base}\\
g_s^{\mathrm{qry}}(\Delta g_s,x)
&=g_s^{(0)}+\Delta g_s
+\mu_s^{\mathrm{cal}}(x-\tau_s).
\label{eq:g_query}
\end{align}
式中，$g_s^{\mathrm{base}}(x)$ 为低速标定路径上的基准间隙，$g_s^{\mathrm{qry}}(\Delta g_s,x)$ 为给定高速间隙增量后的查询间隙，$\Delta g_s$ 为第 $s$ 个传感器处高速状态相对于低速标定状态的等效径向间隙变化。两式使用同一个 $\mu_s^{\mathrm{cal}}$，因此高速径向间隙变化表示为低速标定路径沿 $g$ 方向的整体平移，同时保留径向间隙随局部周向位置变化的关系。

在相同空间查询位置 $x$ 处，将高速查询间隙与低速基准间隙对应的响应面输出相减，定义间隙变化引起的电压增量为
\begin{align}
\Delta\Rcal_s(\Delta g_s,x)
={}&\Rcal\!\left(
g_s^{\mathrm{qry}}(\Delta g_s,x),
x_s^{\mathrm{lib}}(x)
\right)
\nonumber\\
&-\Rcal\!\left(
g_s^{\mathrm{base}}(x),
x_s^{\mathrm{lib}}(x)
\right).
\label{eq:response_increment}
\end{align}
式中，$\Delta\Rcal_s(\Delta g_s,x)$ 表示静态响应面给出的相对电压变化。将该变化经传感器增益映射后叠加到实测低速模板上，定义传感器相关的间隙增量前向模型为
\begin{equation}
\Fsen_s(\Delta g_s,x)
=T_s^{\mathrm{low}}(x)
+\kappa_s\Delta\Rcal_s(\Delta g_s,x),
\label{eq:sensor_model}
\end{equation}
当 $\Delta g_s=0$ 时，式\eqref{eq:response_increment}严格为零，式\eqref{eq:sensor_model}退回实测模板 $T_s^{\mathrm{low}}(x)$。因此，低速模板保存绝对基准波形，静态响应面只提供相对低速状态的间隙增量响应，二者不会重复定义基准波形。

$T_s^{\mathrm{low}}(x)$、$g_s^{(0)}$、$\mu_s^{\mathrm{cal}}$、$\tau_s$、$\zeta_s$ 和 $\kappa_s$ 均在高速识别前确定，并在同一次高速主拟合中保持固定。重复低速标定和参数扰动用于评价这些标定量对高速间隙与振动结果的影响。

\subsection{高速完整波形联合模型}

\subsubsection{周向振动位移表示}

本文将振动识别量定义为叶尖在周向上的等效位移，并把同步与异步成分的相位自变量分开。同步分量随实测累计转角 $\theta_{sn}$ 演化，异步分量随绝对时间 $t_{sn}$ 演化：
\begin{align}
u_k^{\mathrm S}(t_{sn})
&=a_k^{\mathrm S}\sin(k\theta_{sn})+b_k^{\mathrm S}\cos(k\theta_{sn}),\\
u_f^{\mathrm A}(t_{sn})
&=a_f^{\mathrm A}\sin\{2\pi f(t_{sn}-t_c)\}
+b_f^{\mathrm A}\cos\{2\pi f(t_{sn}-t_c)\}.
\label{eq:vibration}
\end{align}
其中 $k\in\mathbb Z_+$ 为离散发动机阶次（engine order，EO），$f$ 为连续异步频率，$t_c$ 为窗口中心。采用实测 $\theta_{sn}$ 可把转速波动纳入同步相位；将同步项写成 $\sin(2\pi k\bar f_rt)$ 只在严格恒速时等价。任一分量的幅值和参考相位由
\begin{equation}
A=\sqrt{a^2+b^2},
\qquad
\varphi=\operatorname{atan2}(b,a)
\label{eq:amp_phase}
\end{equation}
得到。

冻结程序不在一次求解中同时推断任意成分数和任意结构，而是在以下四类预先给定的结构内辨识未知频率参数：
\begin{align}
\mathcal M_{\mathrm{SS1}}(k_1):\quad
u(t)&=u_{k_1}^{\mathrm S}(t),\label{eq:model_single_sync}\\
\mathcal M_{\mathrm{AA1}}(f_1):\quad
u(t)&=u_{f_1}^{\mathrm A}(t),\label{eq:model_single_async}\\
\mathcal M_{\mathrm{SS2}}(k_1,k_2):\quad
u(t)&=u_{k_1}^{\mathrm S}(t)+u_{k_2}^{\mathrm S}(t),
\quad k_1<k_2,\label{eq:model_dual_sync}\\
\mathcal M_{\mathrm{SA}}(k_1,f_2):\quad
u(t)&=u_{k_1}^{\mathrm S}(t)+u_{f_2}^{\mathrm A}(t),
\quad |f_2-k_1\bar f_r|\geq\Delta f_{\min}.
\label{eq:model_sync_async}
\end{align}
单同步和双同步模型在给定频带内枚举整数 EO；单异步模型在连续时间频率网格上生成候选，并在完整模型中释放 $f$ 精修；同步--异步模型同时枚举整数 EO 和异步频率候选。$\Delta f_{\min}$ 仅用于避免同一有限窗口内明显重复的同步/异步表示，其取值由记录长度和分辨性审计决定，不是普适可分辨极限。

四类模型的结构已知意味着只在相应模型族内搜索，并不意味着真实 $k$ 或 $f$ 被提供给算法。跨结构类别的自动选择属于更高一层模型阶数问题，当前正式证据不支持把它与结构内频率辨识混为一个无条件保证的问题。对于双分量模型，还要求估计幅值不低于预定物理容差下限；该约束可排除数值上接近零的伪分量，但不能替代候选子空间分离度判断。

\subsubsection{高速波形观测方程}

按照第2节的建模假设，高速相对于低速的主要静态变化表示为径向间隙增量，叶片振动表示为周向等效位移。由于各传感器的局部坐标已在第3.2节统一，全部传感器共享同一周向振动位移函数 $u(t)$。

叶片振动和空间对准误差通过改变模板取值位置进入前向模型。第 $s,n$ 个观测点的有效周向坐标定义为
\begin{equation}
\xi_{sn}=x_{sn}-d-u(t_{sn}),
\label{eq:xi}
\end{equation}
其中，$d$ 为当前数据窗口内全部传感器共享的空间对准偏移，用于表征并补偿同一坐标构造过程产生的公共零位偏差。各传感器的固定安装角差已通过式\eqref{eq:local_space_coordinate}中的名义参考转角 $\theta_{s\ell}^{\mathrm{ref}}$ 纳入局部坐标定义，因此高速主模型不再引入逐传感器空间偏移。本文规定 $u(t)>0$ 表示叶尖沿局部周向坐标正方向移动；在固定名义观测坐标下，模板与响应面的有效取值位置相应减小，故式\eqref{eq:xi}中振动位移前取负号。

将式\eqref{eq:xi}代入式\eqref{eq:sensor_model}，得到固定低速标定梯度模型的完整波形预测：
\begin{equation}
\widehat v_{sn}(\bm\theta_c;\bm k)
=\Fsen_s\!\left(
\Delta g_s,
x_{sn}-d-u(t_{sn})
\right).
\label{eq:observation}
\end{equation}
式\eqref{eq:observation}是本文的核心前向模型。高速间隙增量 $\Delta g_s$ 改变响应面的间隙输入，公共空间偏移和振动位移改变模板及响应面的空间取值位置，固定的 $\mu_s^{\mathrm{cal}}$ 则在该有效周向位置处参与间隙路径计算。因此，已标定梯度与振动位移之间的几何作用被保留，但梯度本身不作为高速自由参数。上述效应通过同一个非线性前向模型共同形成预测波形，不能解释为彼此独立的加性电压项。

\subsubsection{识别参数与目标函数}

记预先给定的结构类别为 $m\in\{\mathrm{SS1},\mathrm{AA1},\mathrm{SS2},\mathrm{SA}\}$，结构内离散 EO 指标为 $\bm k_m$，其有限候选集合为 $\mathcal K_m$；纯异步模型不含离散 EO 时取 $\mathcal K_m=\{\varnothing\}$。给定 $(m,\bm k_m)$ 后，连续待估参数统一写为
\begin{equation}
\bm\theta_c=
\left[
\Delta g_1,\ldots,\Delta g_{N_s},d,
\bm c_m^{\mathrm T},\bm f_{m,\mathrm A}^{\mathrm T}
\right]^{\mathrm T}.
\label{eq:theta}
\end{equation}
其中 $\bm c_m$ 收集该结构全部正弦/余弦系数，$\bm f_{m,\mathrm A}$ 收集异步连续频率；参数中不包含高速自由梯度增量。记具有有效间隙标定信息的传感器集合为 $\mathcal S_g$；对 $s\notin\mathcal S_g$，在整个高速主优化中约束 $\Delta g_s\equiv0$，但相应观测仍参与公共偏移和振动参数估计。为区分候选生成和最终参数估计，变量投影阶段可使用预先确定的权重矩阵减小低灵敏度点的数值影响；最终候选比较采用同一可信样本上的普通完整电压残差，并加入支撑域可行性约束和预先规定的正则项。定义
\begin{equation}
\Phi(\bm\theta_c;\bm k)=
\sum_{s,n}
\left[v_{sn}-\widehat v_{sn}(\bm\theta_c;\bm k)\right]^2
+P_{\mathrm{out}}(\bm\theta_c;\bm k)
+P_{\mathrm{invalid}}(\bm\theta_c;\bm k)
+P_g(\Delta\bm g).
\label{eq:objective}
\end{equation}
$P_{\mathrm{out}}$ 用于抑制查询点超出静态响应面和低速模板的有效范围，$P_{\mathrm{invalid}}$ 用于惩罚无效插值，$P_g$ 用于限制缺少独立间隙真值时的过度静态修正。令 $o_{sn}$ 为第 $s,n$ 个查询点到有效标定域的越界距离，$N_{\mathrm{invalid}}$ 为无效查询点数，$N$ 为当前窗口观测点数，可采用
\begin{equation}
\begin{aligned}
P_{\mathrm{out}}&=\gamma_{\mathrm{out}}^2\sum_{s,n}o_{sn}^2,\\
P_{\mathrm{invalid}}&=M_{\mathrm{invalid}}N_{\mathrm{invalid}},\\
P_g&=
\frac{N\gamma_g^2}{N_s}
\sum_{s=1}^{N_s}
\left(\frac{\Delta g_s}{g_{s,\mathrm{scale}}}\right)^2,
\end{aligned}
\label{eq:penalties}
\end{equation}
其中，间隙坐标和周向坐标均换算为长度单位后，$o_{sn}$ 取查询间隙和周向位置相对于各自有效区间的二维欧氏越界量；$\gamma_{\mathrm{out}}$ 的量纲用于将该越界量换算为电压残差尺度。$M_{\mathrm{invalid}}$ 的量纲为电压平方，$\gamma_g$ 为间隙增量正则项的电压尺度；这些系数应在全部数据集上保持固定。$g_{s,\mathrm{scale}}$ 为由标定范围或独立物理信息给出的间隙尺度。不同方法比较时应保持相同观测点、完整波形残差定义和惩罚口径。式\eqref{eq:objective}中的未加权残差是最终结果口径，不与变量投影阶段的加权残差混用。该口径默认各传感器的有效采样点数和电压尺度具有可比性；否则，总残差会自然偏向点数较多或电压尺度较大的传感器。本文保留未加权总残差作为统一优化准则，同时通过逐传感器误差和传感器组合敏感性检查结果是否受单个传感器主导。

为明确模型层级，固定模板模型对应 $\Delta g_s=0$；固定低速标定梯度模型允许 $\Delta g_s$ 变化但固定 $\mu_s^{\mathrm{cal}}$；允许高速自由梯度增量的模型仅作为敏感性对照，不作为主振动幅值模型。

\subsection{投影局部搜索与完整模型优化}
\label{sec:candidate_solution}

直接从任意初值同时搜索阶次、连续频率、间隙、公共空间偏移和振动参数容易进入局部极小。本文不设置忽略振动的独立静态初值步骤，而是对每组候选阶次及连续频率参数先计算振动候选，再在该候选下构造间隙搜索范围，最后对全部保留候选进行完整非线性优化。

\subsubsection{加权变量投影候选生成}

对给定候选阶次向量 $\bm k$、连续频率偏移向量 $\Delta\bm f$、参考间隙增量 $\Delta\bm g^{\mathrm{ref}}$ 和参考公共空间偏移 $d^{\mathrm{ref}}$，先在参考状态附近对振动位移作一阶展开。$d^{\mathrm{ref}}$ 在预设的有界公共偏移区间内作为外层候选参数；本文取 $\Delta\bm g^{\mathrm{ref}}=\bm 0$，该零值只是投影中心构造所采用的参考状态，不是高速间隙不变的物理假设。参考状态的无振动预测值及其总空间导数定义为
\begin{equation}
F_{sn}^{(0)}
=\Fsen_s(\Delta g_s^{\mathrm{ref}},x_{sn}-d^{\mathrm{ref}}),
\qquad
F_{x,sn}^{(0)}=
\left.\frac{\partial\Fsen_s(\Delta g_s,x)}{\partial x}
\right|_{\substack{\Delta g_s=\Delta g_s^{\mathrm{ref}}\\
x=x_{sn}-d^{\mathrm{ref}}}},
\label{eq:local_derivative}
\end{equation}
式中，$F_{sn}^{(0)}$ 为参考间隙和参考公共偏移下的无振动预测电压，$F_{x,sn}^{(0)}$ 为固定低速标定梯度前向模型对空间查询坐标的总导数。由此，观测点 $s,n$ 处的波形关于周向振动位移的一阶近似为
\begin{equation}
\widehat v_{sn}
\approx F_{sn}^{(0)}-F_{x,sn}^{(0)}u(t_{sn}).
\label{eq:linearized_wave}
\end{equation}
$F_{x,sn}^{(0)}$ 包含低速模板形状、坐标映射以及已标定梯度对空间查询的共同影响。

对给定阶次向量 $\bm k$ 和频率偏移向量 $\Delta\bm f$，式\eqref{eq:linearized_wave}对 $a_q,b_q$ 为线性关系。将全部观测点按固定顺序堆叠，得到线性候选模型
\begin{equation}
\bm r_0(d^{\mathrm{ref}};\Delta\bm g^{\mathrm{ref}})
\approx
\bm D(\bm k,\Delta\bm f,d^{\mathrm{ref}};\Delta\bm g^{\mathrm{ref}})
\bm\alpha+\bm e,
\qquad
\bm\alpha=
[a_1,b_1,\ldots,a_Q,b_Q]^{\mathrm T},
\label{eq:vp_model}
\end{equation}
式中，$\bm v$ 和 $\bm F^{(0)}$ 分别为实测电压与参考预测电压的堆叠向量，$\bm r_0=\bm v-\bm F^{(0)}$ 为参考残差，$\bm\alpha$ 为全部振动线性系数组成的向量，$\bm D$ 为相应设计矩阵，$\bm e$ 为局部线性化余项。矩阵 $\bm D$ 中对应第 $q$ 个正弦和余弦系数的两列定义为
\begin{equation}
D_{sn,2q-1}=-F_{x,sn}^{(0)}\sin\psi_{q,sn},
\qquad
D_{sn,2q}=-F_{x,sn}^{(0)}\cos\psi_{q,sn}.
\label{eq:vp_design}
\end{equation}
变量投影阶段使用预先确定的对角权重矩阵 $\bm\Lambda_{\mathrm{VP}}$，降低低空间灵敏度点和响应面边缘点对候选生成的影响。为使投影和范数均在普通欧氏空间中解释，白化残差和设计矩阵定义为
\begin{equation}
\widetilde{\bm r}_0=\bm\Lambda_{\mathrm{VP}}^{1/2}\bm r_0,
\qquad
\widetilde{\bm D}=\bm\Lambda_{\mathrm{VP}}^{1/2}\bm D.
\label{eq:vp_whitening}
\end{equation}
在每组候选阶次、连续频率偏移和公共空间偏移参考值下，振动线性系数由变量投影求得\cite{Golub1973,Golub2003}：
\begin{equation}
\widehat{\bm\alpha}
(\bm k,\Delta\bm f,d^{\mathrm{ref}};\Delta\bm g^{\mathrm{ref}})
=\widetilde{\bm D}
(\bm k,\Delta\bm f,d^{\mathrm{ref}};\Delta\bm g^{\mathrm{ref}})^{\dagger}
\widetilde{\bm r}_0
(d^{\mathrm{ref}};\Delta\bm g^{\mathrm{ref}})
\label{eq:vp_solution}
\end{equation}
式中，$(\cdot)^{\dagger}$ 表示通过奇异值分解计算的 Moore--Penrose 广义逆。将线性系数代回白化局部模型后，变量投影候选得分定义为
\begin{equation}
\Psi_{\mathrm{VP}}
(\bm k,\Delta\bm f,d^{\mathrm{ref}};\Delta\bm g^{\mathrm{ref}})
=\left\lVert
\widetilde{\bm r}_0-
\widetilde{\bm D}\widehat{\bm\alpha}
\right\rVert_2^2.
\label{eq:vp_score}
\end{equation}
外层按式\eqref{eq:vp_score}比较候选阶次、连续频率偏移和公共空间偏移参考值，内层直接求解 $a_q,b_q$，从而减少候选生成阶段的非线性搜索维数。为保证线性振动系数在当前候选下可数值分离，$\widetilde{\bm D}$ 经统一奇异值阈值截断后的数值秩应为 $2Q$，且最小保留奇异值或条件数满足预设稳定性要求；不满足该条件的组合不进入保留候选集合 $\mathcal C(\bm k)$。

式\eqref{eq:vp_solution}只用于为每组外层候选求得相应的振动线性系数；与该组系数配对的 $d^{\mathrm{ref}}$ 记为 $d^{\mathrm{cand}}$。在满足数值秩条件的组合中，保留投影得分较小的有限候选。具体保留数量或得分阈值属于统一设置的求解参数，不改变后续完整模型的判定准则。这些量均为完整优化的候选初值，不作为最终振动或空间偏移结果。所有保留候选仍需进入式\eqref{eq:observation}的完整非线性模型，不能用一阶近似残差替代最终完整波形残差。

\subsubsection{条件间隙灵敏度与局部搜索范围}

对每个保留候选，将变量投影得到的振动系数与对应的公共空间偏移、频率参数和参考间隙增量组合为
\begin{equation}
\bm\theta^{\mathrm{cand}}=
\left\{
\Delta\bm g^{\mathrm{ref}},
d^{\mathrm{cand}},
\bm k^{\mathrm{cand}},
\widehat{\bm a},
\widehat{\bm b},
\Delta\bm f^{\mathrm{cand}}
\right\},
\label{eq:candidate_point}
\end{equation}
其中，$\widehat{\bm a}$ 和 $\widehat{\bm b}$ 由式\eqref{eq:vp_solution}中的 $\widehat{\bm\alpha}$ 分解得到。在该候选点代入完整非线性前向模型，定义
\begin{equation}
\bm r^{\mathrm{cand}}
=\bm v-\widehat{\bm v}(\bm\theta^{\mathrm{cand}}).
\label{eq:candidate_residual}
\end{equation}
下文的 $\bm J_g$ 和 $\bm J_{\nu}$ 均定义为预测波形 $\widehat{\bm v}$ 对相应参数的雅可比，并在 $\bm\theta^{\mathrm{cand}}$ 处由完整模型求导或差分得到，即
\begin{equation}
\bm J_g=
\left.
\frac{\partial\widehat{\bm v}}
{\partial\Delta\bm g}
\right|_{\bm\theta^{\mathrm{cand}}},
\qquad
\bm J_{\nu}=
\left.
\frac{\partial\widehat{\bm v}}
{\partial\bm\nu}
\right|_{\bm\theta^{\mathrm{cand}}}.
\label{eq:prediction_jacobians}
\end{equation}
离散阶次 $\bm k^{\mathrm{cand}}$ 在局部投影中保持固定，不构造对阶次的微分灵敏度。为简化记号，以下将 $\bm r^{\mathrm{cand}}$ 记为 $\bm r$。当参数修正量为 $\delta\bm g$ 和 $\delta\bm\nu$ 时，修正后的残差在该候选点附近满足
\begin{equation}
\bm r_{\mathrm{new}}
\approx
\bm r-
\bm J_g\delta\bm g-
\bm J_{\nu}\delta\bm\nu,
\label{eq:local_model}
\end{equation}
其中，$\bm J_g$ 由具有间隙标定信息的传感器集合 $\mathcal S_g$ 对应的间隙灵敏度列组成。干扰参数灵敏度矩阵写为
\begin{equation}
\bm J_{\nu}=
[\bm J_d,\bm J_{a},\bm J_b,\bm J_{\Delta f}],
\label{eq:nuisance_jacobian}
\end{equation}
其各列分别对应公共空间偏移、全部振动正弦和余弦系数以及连续频率偏移。$\delta\bm g$ 和 $\delta\bm\nu$ 为参考状态附近的局部参数修正量。未进入 $\mathcal S_g$ 的传感器仍参与公共偏移、振动和频率方向的估计，但其间隙增量固定为零。

采用与变量投影相同的权重进行白化：
\begin{equation}
\widetilde{\bm r}=\bm\Lambda_{\mathrm{VP}}^{1/2}\bm r,
\qquad
\widetilde{\bm J}_{g}=\bm\Lambda_{\mathrm{VP}}^{1/2}\bm J_g,
\qquad
\widetilde{\bm J}_{\nu}=\bm\Lambda_{\mathrm{VP}}^{1/2}\bm J_{\nu}.
\label{eq:gap_whitening}
\end{equation}
通过奇异值分解得到 $\widetilde{\bm J}_{\nu}$ 的数值列空间正交基 $\bm U_{\nu}$，并定义
\begin{equation}
\bm P_{\nu}=\bm U_{\nu}\bm U_{\nu}^{\mathrm T},
\qquad
\bm Q_{\nu}=\bm I-\bm P_{\nu}.
\label{eq:projector}
\end{equation}
$\bm P_{\nu}$ 投影到公共偏移、振动和频率灵敏度张成的局部子空间，$\bm Q_{\nu}$ 则保留这些方向不能一阶解释的部分。由此得到联合间隙矩阵和投影残差：
\begin{equation}
\bm G=\bm Q_{\nu}\widetilde{\bm J}_g,
\qquad
\bm r_{\perp}=\bm Q_{\nu}\widetilde{\bm r}.
\label{eq:projected_gap_system}
\end{equation}
对 $\bm G$ 进行奇异值分解，并仅保留大于既定相对阈值的奇异值。记由此得到的截断广义逆为 $\bm G_{\mathrm{tr}}^{\dagger}$。只有当保留的数值秩等于待估间隙参数个数时，才计算联合局部线性间隙修正量
\begin{equation}
\delta\bm g^{\mathrm{LS}}
=\bm G_{\mathrm{tr}}^{\dagger}\bm r_{\perp}.
\label{eq:joint_gap_ls}
\end{equation}
若 $\bm G$ 的数值秩不足，则不计算联合投影搜索中心；对全部 $s\in\mathcal S_g$ 取 $\overline{\Delta g}_s=\Delta g_s^{\mathrm{ref}}$，并采用预设的保守半宽 $h_{g,s}^{\mathrm{cons}}$。奇异值截断阈值在全部数据集上保持一致，避免因候选而改变秩判据。
为评价第 $s$ 个间隙参数，在已去除干扰参数方向的基础上，还需排除其他间隙参数的一阶灵敏度方向。记 $\bm G_{-s}$ 为删去第 $s$ 列后的 $\bm G$，$\bm P_{g,-s}$ 为通过奇异值分解得到的 $\bm G_{-s}$ 数值列空间正交投影，$\bm g_s$ 为 $\bm G$ 的第 $s$ 列，$\widetilde{\bm j}_{g,s}$ 为白化间隙雅可比矩阵 $\widetilde{\bm J}_g$ 的第 $s$ 列；若不存在其他有效间隙参数，则取 $\bm P_{g,-s}=\bm0$。由此定义
\begin{equation}
\bm q_{g,s}=
\left(\bm I-\bm P_{g,-s}\right)\bm g_s,
\qquad
\eta_{g,s}=
\frac{\left\lVert\bm q_{g,s}\right\rVert_2}
{\left\lVert\widetilde{\bm j}_{g,s}\right\rVert_2}.
\label{eq:gap_separability}
\end{equation}
$\eta_{g,s}$ 描述第 $s$ 个间隙参数的一阶灵敏度中，不能被公共空间偏移、振动、频率偏移及其他间隙参数灵敏度方向解释的比例。因此，$\eta_{g,s}$ 是当前候选附近的局部条件可辨识性指标；其值较小表示该参数的一阶灵敏度接近混叠，但它本身不是间隙误差上界。即使原始间隙灵敏度列在不同传感器观测区间上互不重叠，经公共干扰子空间投影后也不再假定其严格正交，因此仍需按联合间隙雅可比矩阵进行条件化处理。

相应的条件投影残差为
\begin{equation}
\bm r_{\perp,s}
=\left(\bm I-\bm P_{g,-s}\right)\bm r_{\perp}.
\label{eq:conditional_gap_residual}
\end{equation}
条件化残差与条件灵敏度之比用于定义局部搜索尺度基准：
\begin{equation}
 h_{g,s}^{\mathrm{proj}}
:=\frac{\left\lVert\bm r_{\perp,s}\right\rVert_2}
{\left\lVert\bm q_{g,s}\right\rVert_2}.
\label{eq:gap_bound}
\end{equation}
该比值反映在当前局部线性模型中，未解释残差相对于独立间隙灵敏度的尺度；它不直接作为间隙估计值。利用联合最小二乘得到的有符号修正量更新局部搜索中心：
\begin{equation}
\overline{\Delta g}_s
=\operatorname{clip}\!\left(
\Delta g_s^{\mathrm{ref}}
+\delta g_s^{\mathrm{LS}},
\Delta g_{s,\min}^{\mathrm{adm}},
\Delta g_{s,\max}^{\mathrm{adm}}
\right),
\label{eq:gap_center}
\end{equation}
其中，$\operatorname{clip}$ 表示按传感器分量投影到各自允许区间；$\Delta g_{s,\min}^{\mathrm{adm}}$ 和 $\Delta g_{s,\max}^{\mathrm{adm}}$ 由响应面覆盖范围与物理允许范围共同确定。本文取 $\Delta g_s^{\mathrm{ref}}=0$，因此投影联合解直接构成候选相关的间隙搜索中心。实际局部搜索半宽设置为
\begin{equation}
h_{g,s}
=\operatorname{clip}\!\left(
\rho_g h_{g,s}^{\mathrm{proj}},
h_{g,\min},
h_{g,\max}
\right),
\label{eq:local_range}
\end{equation}
其中，$\rho_g>1$ 为保守放大系数，$h_{g,\min}$ 和 $h_{g,\max}$ 分别避免搜索范围过窄或超出间隙库支持范围。$h_{g,s}^{\mathrm{proj}}$ 是投影残差驱动的局部搜索尺度，不是统计置信区间，也不是非线性问题中真实间隙误差的严格上界。除上述联合矩阵秩不足时的整体回退外，若仅第 $s$ 个参数的 $\eta_{g,s}$ 或 $\|\bm q_{g,s}\|_2$ 低于阈值，则只对该传感器取 $\overline{\Delta g}_s=\Delta g_s^{\mathrm{ref}}$ 和 $h_{g,s}=h_{g,s}^{\mathrm{cons}}$；其余满足条件的传感器仍采用投影中心和半宽。若干扰子空间对奇异值截断阈值过度敏感，则按联合秩不足处理。若最终结果命中允许边界，也应标记为局部弱可辨识，不能据此声称全局唯一性。

\subsubsection{双同步 Funnel V2 的候选召回职责}

前述变量投影给出通用候选生成思想，但双同步模型的离散 EO 对数量较大，且小位移线性化、固定间隙更新和完整联合更新的排序可能明显不同。冻结仿真采用 Funnel V2，其固定层级写为
\begin{equation}
\begin{aligned}
1496\ \mathrm{VP}
&\rightarrow272\ (\mathrm{exact0+fixed\text{-}g\ GN1})
\rightarrow136\times1\ (\mathrm{joint\ GN})\\
&\rightarrow24\times1\ (\mathrm{joint\ GN})
\rightarrow3\ (\mathrm{full\ refinement}).
\end{aligned}
\label{eq:funnel_v2}
\end{equation}
在当前频带中共有 136 个互异 EO 对和 11 个间隙状态，故初始 VP 组合数为 $136\times11=1496$。每个 EO 对先保留两个相互分离的间隙状态，执行完整前向回放和一次固定间隙 GN；合并间隙状态后，\emph{全部} 136 个 EO 对均执行一次允许间隙、公共偏移、两分量幅值和相位共同更新的 joint--GN。只有在这一步之后才按完整电压残差删减 EO 对。随后 Top--24 再执行一次 joint--GN，Top--3 进入有界完整非线性最小二乘精修。

Funnel V2 解决的是候选召回问题：若真 EO 在固定间隙 GN1 后排名较低，第一次全 EO joint--GN 仍可把它恢复到最终精修集合。令 $\mathcal C_3$ 为 Top--3 集合，则召回事件为
\begin{equation}
\operatorname{Recall@3}
=\mathbb I\{\bm k^{\mathrm{true}}\in\mathcal C_3\}.
\label{eq:recall_at_3}
\end{equation}
即使式\eqref{eq:recall_at_3}成立，最终选择仍由各候选完整精修后的普通电压残差决定；当错误 EO 模型在含噪或含确定性模板偏差的数据上取得更小残差时，扩大 Top--24 或 Top--3 不会改变竞争结果。因此，Funnel V2 不承担“近退化模型必选真值”的职责，也不应通过继续增加候选预算来掩盖信息不足。

\subsubsection{完整模型有界优化与候选选择}

对每个保留的候选，以上述投影中心和半宽构造间隙边界，并在候选附近对完整前向模型进行有界联合优化。主要局部范围写为
\begin{equation}
\begin{aligned}
|\Delta g_s-\overline{\Delta g}_s|&\leq h_{g,s},\\
|d-d^{\mathrm{cand}}|&\leq h_d,\\
|f_{m,\mathrm A,q}-f_{m,\mathrm A,q}^{\mathrm{cand}}|&\leq h_{f,q}.
\end{aligned}
\label{eq:search_domain}
\end{equation}
$d^{\mathrm{cand}}$ 和 $f_{m,\mathrm A,q}^{\mathrm{cand}}$ 分别为候选生成阶段给出的公共偏移和异步连续频率参考值；纯同步模型不含后一项。振动系数、间隙增量、公共偏移和异步频率均在式\eqref{eq:observation}的完整非线性模型中更新。加权残差只用于候选生成和投影搜索范围构造；每个保留候选的局部优化和最终排序均采用式\eqref{eq:objective}规定的同一普通完整电压目标，从而保证候选局部优化准则与最终排序准则一致。

对给定离散阶次向量 $\bm k$，记保留候选集合为 $\mathcal C(\bm k)$，由第 $j$ 个候选及式\eqref{eq:search_domain}构成的局部有界域为 $\Omega_j(\bm k)$。各候选的完整模型局部解及最终候选选择分别写为
\begin{equation}
\widehat{\bm\theta}_{c,j}(\bm k)
=\arg\min_{\bm\theta_c\in\Omega_j(\bm k)}
\Phi(\bm\theta_c;\bm k),
\qquad
(\widehat{\bm k},\widehat j)
=\arg\min_{\substack{\bm k\in\mathcal K\\j\in\mathcal C(\bm k)}}
\Phi\!\left(\widehat{\bm\theta}_{c,j}(\bm k);\bm k\right).
\label{eq:discrete_continuous_objective}
\end{equation}
$\Omega_j(\bm k)$ 同时受总体标定范围、物理边界和候选局部范围约束，最终主参数记为 $\widehat{\bm\theta}=\{\widehat{\bm\theta}_{c,\widehat j}(\widehat{\bm k}),\widehat{\bm k}\}$。式\eqref{eq:discrete_continuous_objective}表示有限候选之间的局部优化与选择，不声称由单一初值获得连续参数问题的全局极小值。

最终输出为
\begin{equation}
\left\{
\widehat{\Delta g}_s,
\widehat d,
\widehat{\bm k}_m,
\widehat{\bm f}_{m,\mathrm A},
\widehat A_q,
\widehat\varphi_q
\right\},
\label{eq:outputs}
\end{equation}
其中同步频率由 $\widehat k_q\bar f_r$ 报告，异步频率由 $\widehat{\bm f}_{m,\mathrm A}$ 报告，$\widehat A_q$ 和 $\widehat\varphi_q$ 由式\eqref{eq:amp_phase}计算。若多个传感器在物理上应共享同一个高速间隙变化，可进一步施加 $\Delta g_s=\Delta g$ 或跨传感器正则化约束；该约束应由传感器布置和被测结构决定，不能仅因其降低残差而采用。

\section{可辨识性分层与可检验条件}
\label{sec:identifiability}

\subsection{四个不同层次}

本文把“能否辨识”拆分为四个互不替代的层次。
\begin{enumerate}
  \item \textbf{结构可辨识性：}在无噪声、前向模型正确且参数位于允许域内时，不同参数或不同离散结构是否产生不同观测。
  \item \textbf{给定离散结构后的连续参数条件可恢复性：}当 EO 或结构类别固定为真值时，间隙、偏移、幅值、相位和异步频率的局部雅可比是否具有足够数值秩及可接受条件数。
  \item \textbf{实用可辨识性：}在有限记录、有限幅值、DAQ 噪声和模板误差存在时，参数误差是否仍小于预定工程阈值。
  \item \textbf{离散 EO 模型选择：}多个经完整优化的 EO 模型之间，真模型是否取得最小最终残差。该层受模型流形距离、确定性模板偏差方向和噪声实现共同影响。
\end{enumerate}
给定真 EO 后连续参数恢复成功，只证明第二层在该工况成立，不能推出盲 EO 选择必然成功；单个随机种子选错 EO 也只说明第三或第四层失败，不能直接推出第一层结构不可辨识。

\subsection{连续参数的局部秩与条件数}

固定结构 $m$ 和离散 EO $\bm k$，令完整预测向量为 $\bm F_{m,\bm k}(\bm\theta_c)$，其在真参数附近的白化雅可比为
\begin{equation}
\widetilde{\bm J}_{m,\bm k}
=\bm\Lambda^{1/2}
\frac{\partial\bm F_{m,\bm k}}{\partial\bm\theta_c}.
\label{eq:full_jacobian}
\end{equation}
若支撑域内 $\bm F$ 连续可微、真参数位于约束域内部，且
\begin{equation}
\operatorname{rank}(\widetilde{\bm J}_{m,\bm k})
=p_m,
\qquad
\sigma_{\min}(\widetilde{\bm J}_{m,\bm k})>0,
\label{eq:local_rank_condition}
\end{equation}
其中 $p_m$ 为连续参数维数，则式\eqref{eq:local_rank_condition}是局部一阶可辨识的可检验充分条件；秩亏是局部一阶不可辨识的必要警报。小扰动下有近似误差界
\begin{equation}
\|\delta\bm\theta_c\|_2
\lesssim
\frac{\|\bm e\|_2}{\sigma_{\min}(\widetilde{\bm J}_{m,\bm k})},
\qquad
\kappa_{\mathrm{loc}}
=\frac{\sigma_{\max}(\widetilde{\bm J}_{m,\bm k})}
{\sigma_{\min}(\widetilde{\bm J}_{m,\bm k})},
\label{eq:local_condition_bound}
\end{equation}
其中 $\bm e$ 汇总 DAQ 噪声、低速模板误差和静态响应面误差。$\kappa_{\mathrm{loc}}$ 大时，即使结构上满秩，幅值、相位或间隙也可能在有限噪声下不稳定，这属于实用退化而非结构不可辨识。

\subsection{候选 EO 子空间分离与模型选择翻转}

为比较两个离散候选 $r$ 和 $c$，先从各自局部雅可比中投影掉共享分量、间隙和公共偏移等共同干扰方向，再分别取数值列空间正交基 $\bm U_r$ 和 $\bm U_c$。二者的最大典型相关系数和最小主角定义为
\begin{equation}
\rho_{\max}(r,c)=\sigma_{\max}(\bm U_r^{\mathrm T}\bm U_c),
\qquad
\vartheta_{\min}(r,c)=\arccos\rho_{\max}(r,c).
\label{eq:principal_angle}
\end{equation}
当 $\rho_{\max}\rightarrow1$ 或 $\vartheta_{\min}\rightarrow0$ 时，两个候选在当前传感器布局与记录窗口下具有近共线敏感度方向。此时离散 EO 仍可能在严格数学意义上不同，但区分它们所需的判别方向能量很小。

设真候选和竞争候选经各自连续参数优化后的无扰动预测之差为 $\bm d_{rc}$，单位判别方向为 $\bm h_{rc}=\bm d_{rc}/\|\bm d_{rc}\|_2$。若低速模板离散化和静态响应近似引入确定性误差 $\bm b$，则其直接影响模型选择的分量为
\begin{equation}
\eta_{rc}=\langle\bm b,\bm h_{rc}\rangle,
\qquad
\rho_b=\frac{|\eta_{rc}|}{\|\bm b\|_2}.
\label{eq:discriminant_projection}
\end{equation}
分箱、SG 与 PCHIP 的总模板 RMSE 即使很小，只要 $\eta_{rc}$ 与真模型的残差优势同量级且方向相反，也可使竞争 EO 的最终残差更低。因而“模板整体误差小”不是模型选择稳定的充分条件；还必须检查误差在 EO 判别方向上的投影。PCHIP 不会消除硬分箱造成的确定性偏差，只会给出连续查询表示。

更一般地，令真参数处观测为 $\bm y=\bm F_r(\bm\theta_r^*)+\bm e$，竞争模型流形到真无扰动观测的距离为
\begin{equation}
D_{rc}=\inf_{\bm\theta_c\in\Omega_c}
\|\bm F_r(\bm\theta_r^*)-\bm F_c(\bm\theta_c)\|_2,
\qquad
D_{\min}=\min_{c\ne r}D_{rc}.
\label{eq:model_manifold_distance}
\end{equation}
在候选域完备且真模型能够精确表示数据时，$D_{\min}>0$ 是无噪声残差唯一选择真结构的基本分离条件；若 $D_{\min}=0$，仅凭当前观测不能唯一选择。由三角不等式可得一个保守充分条件：
\begin{equation}
D_{\min}>2\|\bm e\|_2
\quad\Longrightarrow\quad
\text{真模型的残差严格小于任一竞争模型}.
\label{eq:model_selection_sufficient}
\end{equation}
实际中 $\bm e$ 包含未知确定性偏差，式\eqref{eq:model_selection_sufficient}主要作为审计准则而非可直接宣称的概率保证。可用成对完整精修后的残差裕量、局部判别距离、式\eqref{eq:principal_angle}和式\eqref{eq:discriminant_projection}共同评价风险。

\subsection{成功辨识的必要条件、充分条件与验收判据}

以下条件在当前方法中具有不同逻辑地位。
\begin{enumerate}
  \item \textbf{运行前必要条件：}可信高速坐标域非空；所有允许查询均被低速模板和静态响应面支撑；间隙物理域被静态库覆盖；低速模板重复性和导数质量通过阈值；候选频带覆盖真频率；观测转数不低于预设最小值。
  \item \textbf{连续参数的局部充分条件：}在真离散结构下，式\eqref{eq:local_rank_condition}成立，$\sigma_{\min}$、间隙条件分离度和有效幅噪比均超过预设阈值，优化从保留候选收敛到同一内部解，且未命中支撑或物理边界。
  \item \textbf{离散模型选择的保守充分条件：}真 EO 被 Funnel 召回，所有入围模型均完成同口径精修，并满足式\eqref{eq:model_selection_sufficient}或经重复噪声/独立窗口验证的正残差裕量。该条件较强，当前仿真并未对所有 EO 对验证。
  \item \textbf{可操作验收判据：}最终成功按预先固定的工程阈值逐项判断，而不以 ``ambiguous'' 标签代替：EO 必须正确；异步频率误差不超过 $\varepsilon_f$；幅值、相位和相对间隙误差分别不超过 $\varepsilon_A$、$\varepsilon_\varphi$ 和 $\varepsilon_g$；同时要求优化收敛、查询域有效且无不可接受的边界命中。
\end{enumerate}
形式化地，仿真成功事件定义为
\begin{equation}
\begin{aligned}
\mathcal E_{\mathrm{succ}}
={}&\mathcal E_{\mathrm{EO}}
\cap\{\|\widehat{\bm f}-\bm f^*\|_\infty\leq\varepsilon_f\}
\cap\{\|\widehat{\bm A}-\bm A^*\|_\infty\leq\varepsilon_A\}\\
&\cap\{\|\operatorname{wrap}(\widehat{\bm\varphi}-\bm\varphi^*)\|_\infty\leq\varepsilon_\varphi\}
\cap\{|\widehat{\Delta g}-\Delta g^*|\leq\varepsilon_g\}
\cap\mathcal E_{\mathrm{num}}.
\end{aligned}
\label{eq:success_event}
\end{equation}
``ambiguous'' 或噪声归一化残差裕量只作为近竞争诊断量。现有残差裕量启发式已被证明会把大量参数正确案例也标为模糊，故不作为主要失败判据，更不能覆盖式\eqref{eq:success_event}的逐参数验收。

\section{仿真协议与正式证据}

正式仿真使用完整 DAQ 时间记录加噪协议。低速记录包含 20 个完整转子周期，高速记录包含 8 个完整转子周期；低速和高速分别按各自完整无噪声记录的交流有效值设置独立高斯噪声。噪声加入后才执行 OPR 映射、支撑域认证和高速可信窗口截取。主结果使用现有三传感器布局、全部可信样本、10、15 和 20 dB 名义 DAQ 信噪比，以及不低于 $0.1\,\mathrm{mm}$ 的振动分量幅值。

在上述主范围内，单同步、单异步、双同步和同步--异步各 9 组代表性测试，以及 6 组弱双分量补充测试，共 42 组均满足式\eqref{eq:success_event}规定的参数判据；最大频率误差为 $0.0935\,\mathrm{Hz}$，最大间隙绝对误差约为 $6.2\times10^{-4}\,\mathrm{mm}$。该证据支持“在所列三传感器、8 转、10--20 dB 和幅值下限范围内成功”，不支持“任意频率、任意布局和任意噪声均唯一可辨识”。5 dB 边界、额外五传感器布局和更长记录只作为适用边界研究。

低速模板基准测试比较硬分箱均值、自适应 SG、局部二次和样条方法。在 30 个组合上，自适应 SG 的平均模板 RMSE 为 $1.33\times10^{-4}\,\mathrm V$，平均导数 RMSE 为 $2.68\times10^{-3}\,\mathrm{V/mm}$，平均导数相关系数为 $0.9988$；这些结果支持自适应 SG 作为冻结低速模板方法。它们不构成任意窗口、任意 OPR 抖动下的无偏证明。

双同步 Funnel V2 的消融表明：第一次全 EO joint--GN 可恢复在固定间隙 GN1 后排名 45--58 的真 EO，从而消除早期 Top--K 截断造成的候选漏检。在 15 dB、20 个噪声种子的代表性弱振幅案例中，真 EO 均进入 Top--3 并正确恢复；在最终竞争审计的困难案例中，真 EO 也基本保持 Top--3 召回。因此，现有证据支持 Funnel V2 的候选召回职责，但不支持其对最终近退化残差竞争提供保证。

X01 与 X02 用于区分两种失效机制。X01 中，真 $(10,12)$ 与竞争 $(10,18)$ EO 的条件典型相关系数可达 $0.9980$--$0.99997$，最小主角约 $0.46^\circ$--$3.60^\circ$。硬分箱 $0.02\,\mathrm{mm}$ 后，即使无 DAQ 噪声，模板确定性偏差也使错误 EO 的完整残差低约 $1.43\times10^{-5}\,\mathrm V$；把低速映射窗口改为支撑感知设置后，真 EO 恢复。这支持“模板离散化偏差与近共线判别子空间共同导致模型选择翻转”，不能简化为低 SNR 不可辨识。当前合成 OPR 映射直接使用真 OPR 时刻，故该审计尚不能支持“OPR 抖动导致翻转”的结论。

X02 的真 $(15,18)$ 与竞争模型同样具有高子空间相关性，但在无噪声 oracle 比较中真模型仍保有正残差距离；固定真 EO 后，25 和 15 dB 下 20/20 连续参数恢复成功，5 dB 下为 14/20。故 X02 的主要边界是弱分量下的有限噪声模型选择和连续参数实用条件，而不是已证明的结构不可辨识。X01 真 EO 固定后的参数恢复也显著改善，进一步说明“离散结构选择失败”和“连续参数模型失效”必须分开报告。

最终结果表同时报告真 EO 在固定间隙 GN1 后的排名、第一次 joint--GN 后的排名、Recall@3、最终 EO、频率/幅值/相位/间隙误差和前两名完整残差裕量。残差裕量用于解释竞争强度，不取代式\eqref{eq:success_event}。

\section{试验验证方法}

试验标定分为两个步骤。第一步，由满足低振动和重复性条件的低速旋转数据经公共空间网格重采样和多转聚合得到实测模板 $T_s^{\mathrm{low}}(x)$；第二步，将该模板与预先建立的静态多间隙响应面按式\eqref{eq:low_calibration}进行有界匹配，得到 $g_s^{(0)}$、$\mu_s^{\mathrm{cal}}$、$\tau_s$、$\zeta_s$ 和 $\kappa_s$。因此，上述标定参数由低速模板和静态间隙库共同确定。试验应明确报告所采用的聚合算子、异常转次剔除规则、有效低速转次数、模板重复性、标定边界、不同初值的一致性和重复标定误差；高速识别阶段保持 $\mu_s^{\mathrm{cal}}$ 固定。

主对比应包括固定模板模型和固定低速标定梯度模型。报告量包括识别阶次、连续频率、振动幅值、总未加权完整波形残差、$\widehat{\Delta g}_s$、$\eta_{g,s}$、$h_{g,s}^{\mathrm{proj}}$、$h_{g,s}$、边界命中情况以及结果随数据窗口的重复性。为揭示采样点数量和电压尺度差异对总残差构成的影响，并避免总残差指标掩盖单个传感器的拟合差异，还应逐传感器报告
\begin{equation}
\mathrm{RMSE}_s=
\sqrt{\frac{1}{N_{\mathrm{obs},s}}\sum_{n=1}^{N_{\mathrm{obs},s}}
\left(v_{sn}-\widehat v_{sn}\right)^2},
\label{eq:sensor_rmse}
\end{equation}
其中 $N_{\mathrm{obs},s}$ 为第 $s$ 个传感器纳入拟合的有效观测点数。若不同传感器组合得到不同振动幅值，还应分别检查间隙增量、低速等效标定误差和单传感器残差对结果的影响。

若有应变片、激光测振或其他独立振动参考，其作用是验证频率、阶次或幅值。若没有独立间隙真值，则 $\widehat{\Delta g}_s$ 解释为相对于低速基准的模型等效径向间隙变化，不据此评价绝对间隙测量精度。

\section{讨论}

低速模板增量模型的主要作用是把高速识别前可获得的信息与需要由高速波形估计的信息分开。实测低速模板定义零增量基准，静态多间隙响应面只提供相对于该基准的变化量，因此避免用解析或拟合静态模板替换实际低速脉冲。其代价是低速空间离散化误差会固定地进入所有高速候选，故模板函数误差、导数误差和 EO 判别方向投影均是模型的一部分，而不是可忽略的预处理细节。

支撑域条件是前向模型成立的必要条件。扩大低速采集窗口可以增加可认证的高速名义区间，但最终可用区间还要扣除公共偏移、所有振动幅值上界和数值边界裕量。当前 X01 审计表明，采集窗口与硬分箱的组合会改变确定性模板偏差；但由于合成数据使用真 OPR 时刻，尚不能把该现象归因于 OPR 计时误差。真实 OPR 抖动、插值误差和触发同步误差需要单独注入并验证。

四类结构化模型通过先验给定“同步/异步及分量数”减少了模型空间，但没有提供真实 EO 或真实频率。该限制应在论文范围中明确。若结构类别也未知，应在四个模型族之间增加复杂度惩罚、交叉窗口验证或独立物理先验；现有 42 组结果不能直接支持该扩展。

传感器数量、角位置和记录长度共同决定候选频率子空间的独立信息。已有布局审计显示，增加同一布局下的优化次数不能解决所有弱相邻 EO 竞争，而按候选子空间互相关设计非均匀传感器布局能够改善边界案例。这一观察与式\eqref{eq:principal_angle}一致，但目前只得到有限案例的数值证据，尚不是最优布局的普适定理。

最终残差最小模型是数据与冻结前向模型下的选择结果，不等于真物理模型的逻辑证明。近竞争残差应与 Recall@3、principal angles、判别方向偏差投影和给定 EO 的 oracle 恢复共同解释。``ambiguous'' 可提示需要更多观测或外部 EO 先验，但现有噪声归一化裕量启发式不能作为主要失败判据。

\section{结论}

本文形成了与冻结联合辨识程序一致的理论链条：完整 DAQ 记录先加噪，OPR 再把时间样本映射为空间坐标；低速侧在观测支撑内采用完整转次分组交叉验证的自适应 SG 和 PCHIP 连续模板；高速侧以实测低速模板为零增量基准、以静态响应面提供间隙增量，并在单同步、单异步、双同步或同步--异步结构内估计相对间隙与振动参数。

双同步 Funnel V2 通过“全部 EO 对第一次 joint--GN、Top--24 第二次 joint--GN、Top--3 完整精修”提高真 EO 候选召回率，但最终离散模型仍由同口径完整残差竞争。结构可辨识性、给定 EO 的连续参数条件恢复、有限噪声下的实用可辨识性和盲 EO 选择必须分层评价。支撑域覆盖、模板导数质量、有效幅噪比、传感器独立信息、候选子空间主角、局部最小奇异值、优化收敛和残差裕量构成可执行的审计量。

现有仿真支持限定范围内四类结构化模型的参数恢复，也支持 Funnel V2 的召回改进；它同时显示，确定性低速模板偏差在近共线 EO 判别方向上的投影可造成无噪声模型选择翻转。故 X01/X02 不能统一归因于低 SNR，也不能把实用退化表述为结构不可辨识。最终成功应按 EO/频率、幅值、相位和相对间隙的预定误差阈值判断，近竞争残差只作为边界解释。缺少独立间隙真值时，间隙输出仍应表述为相对于低速标定状态的模型等效变化量。

\section*{主要符号}

\begin{longtable}{>{\centering\arraybackslash}p{0.24\textwidth}p{0.68\textwidth}}
\toprule
符号 & 含义\\
\midrule
\endfirsthead
\toprule
符号 & 含义\\
\midrule
\endhead
$s,n,q$ & 传感器、观测点和振动成分编号\\
$N_s,Q$ & 参与识别的传感器数和预设振动成分数\\
$L_s,N_{\mathrm{obs},s}$ & 第 $s$ 个传感器的有效低速转次数和高速有效观测点数\\
$v_{s\ell}^{\mathrm{low}}(t),\overline v_{s\ell}^{\mathrm{low}}(x)$ & 第 $s$ 个传感器第 $\ell$ 个低速转次的时间域原始波形和公共空间网格重采样波形\\
$v_{sn},\widehat v_{sn}$ & 实测电压和完整模型预测电压\\
$s_{\mathrm{DAQ}},\sigma_{\mathrm{DAQ}},\mathrm{SNR}_{\mathrm{DAQ}}$ & 完整 DAQ 无噪声记录的交流有效值、所加噪声标准差和名义采集层信噪比\\
$V_{\mathrm{tip}},\Delta t_u,\delta t_g$ & 叶尖周向速度、振动引起的到达时刻偏差和间隙变化引起的特征定位偏差\\
$t_m^{\mathrm{OPR}},\theta(t),\Delta\theta_m^{\mathrm{OPR}}$ & OPR参考时刻、累计转角和相邻参考位置之间的已知转角增量\\
$\theta_{s\ell}^{\mathrm{ref}},R_{\mathrm{tip}},x_{s\ell n}$ & 目标叶片经过的名义参考转角、叶尖参考半径和局部周向空间坐标\\
$t_{sn},x_{sn},\xi_{sn}$ & 采样时刻、名义局部周向坐标和振动修正后的有效周向坐标\\
$\mathcal X_{L,s},\mathcal X_R,\mathcal X_{H,s}^{\mathrm{cert}}$ & 低速观测支撑、静态响应面空间支撑和高速认证辨识区间\\
$g,x$ & 等效径向间隙和局部周向坐标\\
$T_x(x),S_g(x)$ & 问题分析中低速模板的空间导数和基准间隙路径附近的间隙灵敏度\\
$\phi_j(g)$ & 静态响应面的间隙方向基函数\\
$\Rcal(g,x)$ & 径向间隙--周向位置静态标定响应面\\
$\Fsen_s(\Delta g_s,x)$ & 固定低速标定梯度的传感器相关前向模型\\
$\Gamma_s(x)$ & 低速模板在静态响应面中的等效径向间隙--周向位置查询路径\\
$T_s^{\mathrm{low}}(x)$ & 第 $s$ 个传感器由低速重复波形聚合得到的实测基准模板\\
$g_s^{(0)}$ & 第 $s$ 个传感器的低速等效基准径向间隙\\
$\Delta g_s$ & 高速状态相对低速基准的径向间隙增量\\
$\mu_s^{\mathrm{cal}}$ & 径向间隙随局部周向位置变化的低速等效标定梯度\\
$\tau_s,\zeta_s,\kappa_s$ & 模板--间隙库中心配准、空间尺度修正和电压增益修正参数\\
$b_s,\Omega_s,P_{\mathrm{cal,out}}$ & 低速标定辅助偏置、有界参数域和标定越界惩罚\\
$d$ & 当前数据窗口内全部传感器共享的空间对准偏移\\
$u(t)$ & 叶尖振动在周向上的等效位移\\
$m,\bm k_m,\bm f_{m,\mathrm A}$ & 频率结构类别、结构内同步 EO 向量和异步连续频率向量\\
$\Delta f_{\min}$ & 多频成分在当前分析条件下允许独立报告的最低频率间隔\\
$a_q,b_q$ & 振动正弦、余弦系数\\
$A_q,\varphi_q$ & 振动幅值和参考相位\\
$\lambda_{sn}^{\mathrm{VP}},\bm\Lambda_{\mathrm{VP}}$ & 变量投影阶段的观测点权重和对角权重矩阵\\
$\Psi_{\mathrm{VP}},\mathcal C(\bm k),\Omega_j(\bm k)$ & 变量投影候选得分、保留候选集合和候选局部有界域\\
$\Delta\bm g^{\mathrm{ref}},d^{\mathrm{ref}}$ & 候选生成和局部投影采用的参考间隙增量与公共偏移\\
$\bm\theta^{\mathrm{cand}},\bm r^{\mathrm{cand}}$ & 投影搜索采用的完整模型候选点及其电压残差\\
$\widetilde{\bm r},\widetilde{\bm J}$ & 经 $\bm\Lambda_{\mathrm{VP}}^{1/2}$ 变换后的白化残差和灵敏度矩阵\\
$\mathcal S_g$ & 具有有效间隙标定信息、允许估计间隙增量的传感器集合\\
$\bm J_g,\bm J_{\nu}$ & 预测波形对间隙参数和空间/振动干扰参数的雅可比矩阵\\
$\eta_{g,s}$ & 第 $s$ 个间隙参数的局部条件可辨识性指标\\
$\rho_{\max},\vartheta_{\min},D_{\min}$ & 候选子空间最大典型相关、最小主角和最小模型流形距离\\
$\sigma_{\min},\kappa_{\mathrm{loc}}$ & 完整连续参数雅可比的最小奇异值和局部条件数\\
$\delta g_s^{\mathrm{LS}}$ & 联合投影系统得到的第 $s$ 个局部线性间隙修正量\\
$\overline{\Delta g}_s$ & 投影修正后的局部间隙增量搜索中心\\
$h_{g,s}^{\mathrm{proj}}$ & 由条件投影残差与灵敏度之比得到的局部搜索尺度基准\\
$h_{g,s},h_{g,s}^{\mathrm{cons}}$ & 最终局部优化采用的间隙搜索半宽和预设保守半宽\\
\bottomrule
\end{longtable}

\clearpage
\small
\begin{thebibliography}{99}

\bibitem{Yu2020}
B. Yu, H. Ke, E. Shen, and T. Zhang,
``A review of blade tip clearance-measuring technologies for gas turbine engines,''
\emph{Measurement and Control}, vol. 53, no. 3--4, pp. 339--357, 2020.

\bibitem{Muller1997}
D. Muller, S. Mozumdar, E. Johann, and A. G. Sheard,
``Capacitive measurement of compressor and turbine blade tip to casing running clearance,''
\emph{Journal of Engineering for Gas Turbines and Power}, vol. 119, no. 4, pp. 877--884, 1997.

\bibitem{Golub1973}
G. H. Golub and V. Pereyra,
``The differentiation of pseudo-inverses and nonlinear least squares problems whose variables separate,''
\emph{SIAM Journal on Numerical Analysis}, vol. 10, no. 2, pp. 413--432, 1973.

\bibitem{Golub2003}
G. H. Golub and V. Pereyra,
``Separable nonlinear least squares: The variable projection method and its applications,''
\emph{Inverse Problems}, vol. 19, no. 2, pp. R1--R26, 2003.

\bibitem{Heller2020}
D. Heller, I. A. Sever, and C. W. Schwingshackl,
``A method for multi-harmonic vibration analysis of turbomachinery blades using Blade Tip-Timing and clearance sensor waveforms and optimization techniques,''
\emph{Mechanical Systems and Signal Processing}, vol. 142, Art. no. 106741, 2020.

\bibitem{Diamond2021}
D. H. Diamond, P. S. Heyns, and A. J. Oberholster,
``Constant speed tip deflection determination using the instantaneous phase of blade tip timing data,''
\emph{Mechanical Systems and Signal Processing}, vol. 150, Art. no. 107151, 2021.

\end{thebibliography}

\end{document}
